\section{A Second Target}
\label{sec:target2}

A single audit says something about one algorithm. Whether the instrument
discriminates, and whether the pattern generalises, needs a second target from
a different family. We take the fully informed particle swarm
(FIPS)~\cite{mendes2004fully}, whose published variants distribute a fixed
budget $\phi$ among a particle's neighbours in proportion to a per-neighbour
score. This is the same construction as \eqref{eq:W} in a velocity-based
rather than a diffusion-based optimizer, so the controls transfer unchanged.

We audit \emph{w}FIPS, the published variant that weights each neighbour's
contribution by the quality of its personal best. The topology is
Barab\'{a}si--Albert rather than a regular lattice: on a regular graph every
degree is equal, $\chi_i$ of \eqref{eq:leverage} is identically zero, and a
control would be vacuous by construction rather than by measurement.

Two preconditions were checked before the comparison. The mechanism is active:
the displacement it induces is $\rho=0.22$, against $0.25$--$0.31$ for the
first target, so the two controls neutralise mechanisms of comparable
strength. And \emph{w}FIPS is not distinguishable from plain FIPS in our setup
($\bar{\delta}=+0.043$, $[-0.065,+0.150]$, $p=0.64$), consistent with the
original report that the weighting variants gave no clear gain; this matters
because it is what makes the next result surprising.

\subsection{Large effects that cancel}
\label{sec:target2:result}

At suite level the quality weighting is indistinguishable from its permutation:
$\bar{\delta}=-0.090$ with a $90\%$ interval of $[-0.278,+0.098]$, $p=0.52$,
equivalent to zero only at the small-effect margin $|\delta|<0.33$. Read from
that line alone, the conclusion would be the same as for centrality.

It is not the same. The per-function view (Fig.~\ref{fig:calibration}, fourth
row) shows $13$ of $28$ functions with $|\delta|>0.5$ and $15$ surviving Holm
correction, in \emph{both} directions. Permuting the quality scores is
catastrophic on $F_4$, $F_6$, $F_7$, $F_9$ and $F_{17}$---$\delta=-1.00$,
every paired run, $p=2\times10^{-9}$---and significantly \emph{helpful} on
$F_8$, $F_{21}$, $F_{24}$ ($\delta\approx+0.85$). The effect sorts by problem
class: mean $\delta$ is $-0.431$ on the simple multimodal functions and
$-0.320$ on the hybrids, where permuting hurts, against $+0.314$ on the
compositions and $+0.434$ on the unimodal pair, where it helps.

So the score carries a great deal of information, and the information is
good on some landscape classes and actively misleading on others. Learning
preferentially from the better neighbours helps where the neighbourhood's
quality ordering predicts which direction is worth following, and hurts where
it concentrates the swarm on an incumbent basin---which is what the sign flip
between simple multimodal and composition landscapes describes.

\subsection{Two ways a single view misleads}
\label{sec:target2:views}

The two targets fail the same suite-level test for opposite reasons, and the
contrast is the methodological point.

For centrality, the per-function effects are small everywhere: nothing
happens, anywhere. For quality weighting, the per-function effects are large
almost everywhere and cancel in the mean: a great deal happens, in both
directions. A suite mean alone calls both ``no effect''.

The mirror-image failure appeared in our own adaptive scheme
(Section~\ref{sec:res:calibration}): there the per-function view under Holm
correction credited only $3$ of $28$ wins and read as a weak effect, while the
suite-level estimate was unambiguous ($\bar{\delta}=-0.101$,
$p=1.5\times10^{-5}$, $24$ of $28$ in its favour). Correction across $28$
comparisons is conservative enough to hide a benefit spread thinly over a
suite.

Neither view is sufficient and neither is wrong. A benefit that is diffuse
survives aggregation and not correction; a benefit that is concentrated and
sign-varying survives correction and not aggregation. We report both
throughout and would encourage the same, because each of the two failures is
demonstrated here on real data rather than argued in principle.

\subsection{A synthetic mechanism, for contrast}
\label{sec:target2:synthetic}

The published FIPS family weights by quality or by distance; there is no
degree-weighted member. We constructed one---$w_n=\deg(n)$ on the same
scale-free graph---not as an audit of anyone's claim but to check that the
instrument reports a harmful mechanism as harmful rather than merely as
non-null. It does: $\bar{\delta}=+0.179$ with interval $[+0.074,+0.283]$,
$p=0.014$, and the ablation is better still, uniform weighting beating degree
weighting by $\bar{\delta}=+0.452$ at $p=1.5\times10^{-7}$. Concentrating
influence on the hubs of a scale-free graph accelerates consensus and costs
diversity, which is the mechanism $\chi_i$ measures: degree weighting has the
largest leverage of any scheme we tested ($\chi=0.52$ against $0.40$ for
quality).

Taken together the instrument now separates five outcomes on one suite with
one statistic: inert by construction (a vertex-transitive topology),
active but uninformative (centrality), informative with class-dependent sign
(quality), uniformly harmful (degree), and uniformly beneficial (adaptive
density). A null returned by a test with that range is a statement about the
mechanism.
